\begin{frame}{Анализ рабочих нагрузок LLM в дата‑центре Acme}
	Крупные языковые модели достигли впечатляющих результатов, 
	но их разработка на больших GPU‑кластерах сопряжена с особыми проблемами:
	частые сбои, сложные стратегии параллелизации и несбалансированное использование ресурсов.
	
	В этой работе представлено шестимесячное эмпирическое исследование трассировок из Acme 
	(кластеры Seren и Kalos, всего 4704 GPU A100) и разработаны системные решения 
	для повышения надежности и эффективности.
\end{frame}

\begin{frame}{Контекст и мотивация}
	Большие языковые модели (LLM) требуют миллиардов параметров и огромных объёмов данных, 
	поэтому компании создают крупные GPU‑кластеры для их обучения.
	
	Однако многие выводы из прежних анализов глубокого обучения (DL) неприменимы к LLM из-за:
	\begin{itemize}
		\item новой парадигмы: масштабное самоконтролируемое предварительное обучение (self-supervised pretraining);
		\item унифицированной архитектуры: все LLM используют Transformers;
		\item специализированного программного стека: Deepspeed, Megatron, Alpa и др.
	\end{itemize}
\end{frame}

\begin{frame}{Контекст и мотивация}
	Большие языковые модели (LLM) требуют миллиардов параметров и огромных объёмов данных, 
	поэтому компании создают крупные GPU‑кластеры для их обучения.
	
	Однако многие выводы из прежних анализов глубокого обучения (DL) неприменимы к LLM из-за:
	\begin{itemize}
		\item новой парадигмы: масштабное самоконтролируемое предварительное обучение (self-supervised pretraining);
		\item унифицированной архитектуры: все LLM используют Transformers;
		\item специализированного программного стека: Deepspeed, Megatron, Alpa и др.
	\end{itemize}
\end{frame}